\begin{center}
    \textbf{\uppercase{Abstract}}
\end{center}
\vspace*{1.5cm}


This research analyses the dynamics of Brazilian states following an aggregate uncertainty shock, using two measures of domestic uncertainty and one of global uncertainty. We employ an SVAR model as suggested by \citeonline{baker2016measuring} and \citeonline{barboza2018efeitos}, to obtain the IRF for each state. Subsequently, a multiple linear regression model with cross-sectional data was estimated to relate the states' responses to their respective economic structures. The results indicate that the states react in a heterogeneous manner, regardless of the measure of uncertainty or economic activity. However, domestic uncertainty measures related to market expectations have a greater negative impact on the economy. Moreover, industrial production tends to be more affected when compared to the GDP proxy, the \gls{IBCR}. Finally, states with a higher participation of sectors such as agriculture, extractive industries, manufacturing, financial services, public services, and a greater degree of openness relative to the states' GDP are less impacted. In contrast, states with a higher participation of construction and high levels of indebtedness are more strongly affected.

\vspace*{1.5cm}
\noindent \textbf{Key-words:}  macroeconomics; uncertainty chocks; economic activity; time series, SVAR, business cycle.

\thispagestyle{empty}